% Editable, self-contained LaTeX source. Compile with pdfLaTeX (twice) or Tectonic.
\RequirePackage{fix-cm}
\documentclass[11pt,a4paper]{article}
\usepackage[a4paper,left=27mm,right=27mm,top=24mm,bottom=23mm,headheight=15pt]{geometry}
\usepackage{amsmath,amssymb,array,booktabs,xcolor,fancyhdr,hyperref}
\definecolor{ink}{HTML}{172F3D}
\definecolor{sea}{HTML}{006863}
\definecolor{mist}{HTML}{ECF4F2}
\definecolor{muted}{HTML}{485D63}
\hypersetup{colorlinks=true,urlcolor=sea,linkcolor=sea,pdftitle={How many states fit inside a causal diamond?},pdfauthor={Null-cone research review},pdfsubject={A reader's guide to arXiv:2611.34924v1}}
\pagestyle{fancy}\fancyhf{}
\fancyhead[L]{\sffamily\small\color{muted}GRAVITY ON THE NULL CONE}
\fancyhead[R]{\sffamily\small\color{muted}Research handbook}
\fancyfoot[L]{\sffamily\small\color{muted}26 November 2026\quad /\quad Review of version 1}
\fancyfoot[R]{\sffamily\bfseries\color{sea}\thepage}
\renewcommand{\headrulewidth}{0.35pt}
\setlength{\parindent}{0pt}\setlength{\parskip}{6pt plus 1pt}
\setlength{\abovedisplayskip}{8pt plus 1pt}\setlength{\belowdisplayskip}{8pt plus 1pt}
\raggedbottom\color{ink}
\newcommand{\kicker}[1]{{\sffamily\bfseries\small\color{sea}\MakeUppercase{#1}}\par}
\newcommand{\chapterpage}[3]{\newpage\pdfbookmark[1]{#2}{#1}\kicker{#3}\vspace{2pt}{\sffamily\bfseries\fontsize{23}{27}\selectfont #2\par}\vspace{4pt}}
\newcommand{\subhead}[1]{\par\addvspace{5pt}{\sffamily\bfseries\large\color{sea}#1\par}\nopagebreak\vspace{-3pt}}
\newcommand{\panel}[2]{\par\addvspace{5pt}\begingroup\setlength{\fboxsep}{9pt}\noindent\colorbox{mist}{\parbox{\dimexpr\linewidth-2\fboxsep\relax}{\sffamily\bfseries\color{sea}#1\par\vspace{3pt}\normalfont\color{ink}#2}}\par\endgroup}
\newcommand{\repo}[2]{\href{https://github.com/sprajs/quantum-gravity/blob/main/#1}{#2}}
\newcommand{\source}[1]{\vfill\noindent{\color{muted}\rule{\linewidth}{0.35pt}}\par\vspace{-3pt}{\small\sffamily #1\par}}
\newcommand{\tr}{\operatorname{Tr}}
\newcommand{\rank}{\operatorname{rank}}
\newcommand{\HH}{\mathcal H}
\begin{document}
\pdfbookmark[0]{How many states fit inside a causal diamond?}{start}
\kicker{A reader's guide to an open problem}
{\sffamily\bfseries\fontsize{30}{34}\selectfont How many states fit inside\\a causal diamond?\par}
\vspace{3pt}
{\large A flash of light defines a laboratory. Can the gravity inside it limit how much information the laboratory can hold?\par}

Choose two events on an observer's path. Light leaves the first event and can return by the second. The region accessible between these events is a \emph{causal diamond}. A gravitational wave distorts a ring of neighbouring light rays and slows its expansion. The rays then travel farther to open out to a given area. The paper turns that run from cap to corner into a clock built from the geometry itself.

\begin{center}
\setlength{\unitlength}{0.85mm}
\begin{picture}(161,64)
\color{sea}\linethickness{1.1pt}
\put(51,4){\line(1,1){27}}\put(51,4){\line(-1,1){27}}
\put(24,31){\line(1,1){27}}\put(78,31){\line(-1,1){27}}
\color{muted}\linethickness{0.45pt}
\multiput(25,31)(3,0){18}{\line(1,0){1.5}}
\put(45,10){\line(1,0){12}}\put(45,52){\line(1,0){12}}
\put(12,9){\vector(0,1){44}}
\put(4,55){\makebox(16,4){\sffamily\small time}}
\color{sea}\put(51,4){\circle*{1.8}}\put(51,58){\circle*{1.8}}
\put(78,31){\circle*{1.8}}\put(24,31){\circle*{1.8}}
\color{muted}\put(55,58){\line(1,0){26}}
\put(81,31){\line(1,0){7}}\put(55,4){\line(1,0){26}}
\put(84,58){\makebox(0,0)[l]{\sffamily\small Future tip (final reception)}}
\put(91,31){\makebox(0,0)[l]{\sffamily\small Corner where the cones meet}}
\put(84,4){\makebox(0,0)[l]{\sffamily\small Past tip (initial flash)}}
\end{picture}\par
{\small\color{muted}A spatial slice through the drawing runs horizontally. The two outer points at the widest slice represent a sphere in three spatial dimensions. The short lines near the tips mark the caps.}\par
\end{center}

The question is about \emph{perfectly distinguishable} preparations, represented by orthogonal quantum states. Classically, a maximum clock reading puts an upper limit on the area of the sphere. Counting preparations also requires knowing how many independent states can share an allowed reading.

\panel{What the review finds}{Classical focusing and a particular quantum correlation can be reproduced. The four-dimensional state count still needs a quantum clock that admits physical states and a bound on the number it admits. The review identifies proof gaps, defects in the construction as written and model counterexamples---not a physical refutation of finite-dimensional diamonds. The following pages work through the calculations and testable quantities.}

\source{An AI-assisted technical review of Venkatesa Chandrasekaran, \href{https://arxiv.org/abs/2611.34924v1}{\textsl{An algebraic quantization of gravity from the null cone}}, version 1, 25 November 2026.}

\chapterpage{focusing}{How light measures a region}{Classical geometry}
Label the rays by their launch directions. On a small bundle, write the cross-sectional area per unit solid angle as $r^2$. The function $r(u)$ is its area radius. The parameter $u$ is affine, which means a freely moving light ray obeys the geodesic equation without a parameter-dependent acceleration. A boundary frame fixes how to compare directions and readings.

On a twist-free bundle in four-dimensional vacuum Einstein gravity, the area radius obeys the focusing equation
\[
  \frac{d^2 r}{du^2}=-\frac{r}{2}\,\sigma_{AB}\sigma^{AB}.
\]
The shear $\sigma_{AB}$ measures the change of shape of the bundle at fixed area. Its square is nonnegative. A wave that produces shear therefore reduces the rate at which the rays spread. A round shear-free bundle has $r''=0$.

\subhead{An affine clock}
Cut the expanding light cone from the past tip (the paper's incoming cone) at a cap, with radius $\epsilon r_S$ on this ray, where $0<\epsilon<1$ and $r_S$ is its corner radius. Let $u_c$ denote the affine coordinate at the cap and $u_S$ its value at the corner. Integrating the focusing equation gives
\[
 \underbrace{r'(u_c)(u_S-u_c)}_{u_*}
 =(1-\epsilon)r_S+
 \frac12\int_{u_c}^{u_S}(u_S-w)r(w)\,\sigma_{AB}\sigma^{AB}(w)\,dw .
\]
Changing the units of $u$ rescales the slope and the interval inversely, so their product is unchanged. The integral gives the extra reading caused by shear. For the paper's angularly averaged clock $\bar u_*$ and total corner area $A$, the corresponding classical bound is
\[
 \bar u_*\ \ge\ (1-\epsilon)\sqrt{A/(4\pi)} .
\]
A round shear-free sphere saturates the bound. In a smooth positive-area bundle the integrated identity can survive an area maximum. Using $r/r_S$ as a coordinate along the entire ray requires monotonicity.

\subhead{What the boundary still has to supply}
The first-order bulk equations reduce to Einstein's equations on the branch with an invertible gravitational tetrad. A finite cap brings an additional question. The two-form that determines the brackets of boundary data must be closed. Vary the corner area and two independent radial profiles of a single polarization, as the cap condition permits. The exterior derivative of the displayed two-form (A.14) retains a term at the cap rim. Its value on the polynomial test family is $1/256$ at $\epsilon=1/2$, in the normalization of that calculation. Any repair that retains the moving cap and its joints, or restricts the allowed variations, must also give a generator with the correct boundary flux.

\source{Preprint \S3 and Appendices A--B. \repo{audit/03-2.md}{Focusing derivation} \enspace/\enspace \repo{audit/03-3.md}{Closure test} \enspace/\enspace \repo{audit/cap-crosscheck.md}{Cap variations} \enspace/\enspace \repo{audit/03-4.md}{Hamiltonian and boundary flux}.}

\chapterpage{pairs}{A quantum calculation that closes}{Worked example 1}
Consider a single angular pattern of gravitational radiation and one of its two polarizations. Radial oscillations of this pattern carry integer labels $k=1,2,\ldots$. The example chooses independent Gaussian oscillators in a vacuum state $\Omega$. Write
\[
   s=\frac{A}{4G_N\hbar},\qquad x=\frac1s,\qquad
   v_k=\langle q_k^2\rangle=\frac{x}{2k}.
\]
Here $G_N$ is Newton's constant, $\hbar$ is the reduced Planck constant and the speed of light is one. At a fixed area, decreasing $x$ weakens the coupling. The source $q_k^2-v_k$ subtracts its vacuum mean. Acting on $\Omega$, it creates a pair of quanta in mode $k$ with amplitude $\sqrt{2}\,v_k$. Distinct pairs are orthogonal. The chosen displacement multiplies each by a phase $e^{-2ikx}$.

\subhead{Sum the mode contributions}
With modes through $N$, the overlap of the pair preparation with its displaced copy is
\[
 C_N(x)=2\sum_{k=1}^N v_k^2e^{-2ikx}
       =\frac{x^2}{2}\sum_{k=1}^N\frac{e^{-2ikx}}{k^2}.
\]
Every term has magnitude $x^2/(2k^2)$. Hence the limit $C(x)$ exists, and its error is at most $x^2/(2N)$. In closed form,
\[
 C(x)=\tfrac12 x^2\operatorname{Li}_2(e^{-2ix}),\qquad
 \operatorname{Li}_2(z)=\sum_{k\ge1}\frac{z^k}{k^2}\quad(|z|\le1).
\]
At $x=0.1$ the limiting answer is approximately $0.006703874-0.002609994\,i$. A direct sum through $100{,}000$ radial modes differs from an independent evaluation by $2.50\times10^{-12}$.

\subhead{Why a weak-coupling expansion misleads}
Each phase starts with $1-2ikx$. The coefficient of $x^3$ in $C_N$ is therefore $-iH_N$, where $H_N=1+1/2+\cdots+1/N$. Its magnitude grows without bound as $N$ increases. In the exact answer,
\[
 \operatorname{Im}C(x)=x^3[\log(2x)-1]+O(x^5)\qquad(x\downarrow0).
\]
Modes with $kx$ of order one invalidate an expansion that treats every phase as small. The exact overlap is finite even at zero displacement because the pair weights already fall as $1/k^2$. Holding the displacement fixed independently of $x$ would leave a constant times $x^2$. Testing gravitational self-interactions requires the Einstein vertices and external wave packets in a much larger region.

\source{Preprint \S4.4, equations (4.32)--(4.34). \repo{audit/04-pair.md}{Derivation and error bounds} \enspace/\enspace \repo{checks/pair.py}{Reproduction script} \enspace/\enspace \repo{audit/scattering.md}{Scope of the scattering comparison}.}

\chapterpage{clock}{Who passes the clock test?}{Worked example 2}
In a quantum theory the reading becomes an operator $T$. A maximum allowed reading $\ell$ selects the subspace onto which $P_\ell=\mathbf 1_{[0,\ell]}(T)$ projects. For a normalized preparation $\Omega$, $\|P_\ell\Omega\|^2$ is the probability of passing. The rank of $P_\ell$ counts independent states that pass. These quantities can be calculated separately in a simple model.

\subhead{A corner with one adjustable shape}
Keep the total area $4\pi b^2$ fixed and set the shear to zero. Let a real number $q$ redistribute area over the sphere with density proportional to $e^{q\cos\theta}$, where $\theta$ is the polar angle. Normalize the density to keep the total area fixed. In the model, wavefunctions are square-integrable functions of $q$ with measure $dq$. With $u_0=(1-\epsilon)b$, angular integration gives
\[
    \frac{T(q)}{u_0}=\sqrt{\frac q2\coth(q/2)},\qquad T(0)=u_0 .
\]
The clock acts by multiplication. Set $\ell=1.1u_0$. The accepted interval is $|q|\le q_*\simeq1.6212273$. For a normalized Gaussian with density $e^{-q^2}/\sqrt\pi$, its probability is $\operatorname{erf}(q_*)\simeq0.97814$. Divide that same interval into as many disjoint pieces as desired. The normalized indicator function of each piece is an allowed state and the functions are mutually orthogonal. The exact rank is infinite. In finite matrix approximations, the accepted count roughly doubles each time the matrix size quadruples.

\begin{center}\small
\begin{tabular}{@{}lrrrrr@{}}\toprule
Size of the truncated position matrix & 16 & 64 & 256 & 1,024 & 4,096\\
Accepted eigenvalues & 6 & 12 & 24 & 46 & 94\\\bottomrule
\end{tabular}\end{center}
At $\ell=u_0$ only the point $q=0$ passes, and it has zero measure. Increasing the threshold by ten per cent admits an entire interval. This explains why both the measure and the physical meaning of $q$ matter. To transfer this result to gravity one must exhibit a physical shape observable that survives the constraints, has a continuous component in its spectrum on an allowed state, and has spectral projections commuting with the clock.

\subhead{Can the gravitational clock admit even one state?}
On an expanding light cone the clock integrates the exponential of accumulated squared shear. Squaring a fluctuating field at a point is singular. With selected bare Gaussian modes its mean derivative square grows with radial resolution. A growing mean alone leaves the probability of a low reading undetermined, yet in two regulated quadratic models even the lowest clock eigenvalue diverges with resolution, emptying any fixed window; whether the nonlinear gravitational clock behaves similarly remains open. Vacuum subtraction changes the positivity argument. A constructive answer would exhibit a physical state of finite limiting clock energy, control its overlap with the prepared state and bound its probability at the specified $\ell$.

\source{Preprint \S\S4.3 and 5.1. \repo{audit/regulator-program.md}{Exact shape benchmark} \enspace/\enspace \repo{checks/corner_clock.py}{Eigenvalue counts} \enspace/\enspace \repo{audit/shear-uv.md}{Short-distance shear test}.}

\chapterpage{sewing}{Closing the cut changes the sum}{The state-count argument}
Imagine cutting an experiment in two. To compute the amplitude between specified initial and final states, insert a complete set of intermediate states $|j\rangle$. If $a_j$ and $b_j$ are the components of the two preparations at the cut, then
\[
  \sum_j |a_j^*b_j|\le
     \left(\sum_j|a_j|^2\right)^{1/2}
     \left(\sum_j|b_j|^2\right)^{1/2}.
\]
Each preparation has finite norm. This is why composing the two pieces, called \emph{open sewing}, converges even in an infinite-dimensional space. For instance, composing two identity operators between fixed unit vectors gives their finite inner product. Closing the same identity operator into a loop asks for $\sum_j 1$. The fixed endpoints had supplied an essential restriction.

\subhead{The weight each allowed state carries}
Assume the quantum clock exists, commutes strongly with the area, and enforces the classical bound. Then $0\le s\le s_*$, where $s_*=\pi\ell^2/[(1-\epsilon)^2G_N\hbar]$. Centre a Gaussian filter at $s_b=\pi b^2/(G_N\hbar)$, the value for a round sphere of radius $b$. Any fixed positive width works here. The entropy calculation chooses variance $2s_b$, giving the reflected weight
\[
 W=e^{-s}f_b(s),\qquad
 f_b(s)=\exp[-(s-s_b)^2/(4s_b)].
\]
The minimum $c$ of this weight on the closed bounded interval is strictly positive. For every finite-dimensional projection $F$ in the selected space,
\[
      c\,\rank F\ \le\ \tr(FWF)\ \le\ \rank F.
\]
Every independent state adds at least $c$ to the ordinary, unnormalized trace. Under these hypotheses a finite trace is already equivalent to finite rank. A proof must obtain control of the closed experiment independently.

\subhead{The proposed boost test}
A relative boost changes the reference frame on one side of the cut. The paper predicts a change in a correlation measured across the two sides. If the space were infinite-dimensional, one could group states with very similar area into long cycles. The boost would supply almost a common phase to the dominant states in each partial sum, giving a conflicting prediction. The missing step is to prove that every exhaustive cut order converges to the same independently defined uncut measurement. For the paper's unnormalized signed insertion, bounded partial sums along those cycles already force finite rank without using the sign of the uncut answer.

A comparison space with two area values and infinite multiplicity at both obeys open sewing. Three ways of including every basis state give boost contrasts of approximately $-0.460$, $-0.184$ and $-0.022$. This demonstrates the dependence on summation order in the model. Gravity could impose an additional identity selecting the physical result.

\source{Preprint \S\S4.4 and 5.2. \repo{audit/04-gluing.md}{Open sewing} \enspace/\enspace \repo{audit/dimension-crosscheck.md}{Boost argument} \enspace/\enspace \repo{checks/05-2.py}{Two-area calculation}.}

\chapterpage{entropy}{Which black-hole entropy is counted?}{An independent comparison}
The Bekenstein--Hawking entropy of a large black hole is its horizon area in gravitational units, $s=A/(4G_N\hbar)$, with Boltzmann's constant set to one. If one derived this by counting states of the interior, it would connect the thermodynamic area law to specific microscopic degrees of freedom. The proposal starts from a boost at the corner, which gives states a relative weight $e^{-s}$. Whenever the trace $Z=\tr(e^{-s})$ is finite and nonzero in the selected space, the normalized density matrix and its entropy obey
\[
   \rho=Z^{-1} e^{-s},\qquad
   -\tr(\rho\log\rho)=\langle s\rangle_\rho+\log Z .
\]
The first term follows algebraically. The value and scaling of the second depend on which states the trace includes.

\subhead{One coefficient can be checked directly}
For a four-dimensional Schwarzschild black hole, the local graviton and ghost fluctuations give a curvature coefficient $53/45$ in the convention with $(4\pi)^{-2}$ factored out. Both this calculation and Sen's use that same local heat-kernel coefficient; converting it to the logarithmic entropy agrees with Sen's result for the density of states per unit mass, summed over angular momenta. At fixed reference units the entropy in this ensemble has the large-area form
\[
    S=s+\frac{106}{45}\log s+O(1).
\]
Fixing one angular momentum component at $J_3=0$ gives $167/90$ for the logarithmic coefficient. Keeping only rotational singlets gives $77/90$. These calculable differences specify how spin selection changes the count.

\subhead{Why the area window matters}
The proposed interior trace uses an area packet with Gaussian weight $f_b(s)$. Its standard deviation in $s$ is $\sqrt{2s_b}$. Suppose the boost-weighted density per unit $s$ behaves locally as $Cs^p$ across an interior packet, far from the endpoints. Then
\[
 \int Cs^p e^{-(s-s_b)^2/(4s_b)}\,ds
       \ \sim\ C\sqrt{4\pi}\,s_b^{p+1/2}.
\]
The packet width contributes a factor $\sqrt{s_b}$. For Schwarzschild, $s$ scales as mass squared, so converting a density per unit $s$ into a density per unit mass multiplies by the same square-root factor. Both counts therefore share one logarithmic coefficient, $p+\tfrac12$; no second half is added. Agreement checks that ensemble conversion rather than supplying an independent state count. An independent microscopic calculation must derive the measure for the area mode, the constraints, the ghost and zero modes, and match the interior preparation to the asymptotic mass and spin ensemble.

\source{Preprint \S5.3 and Appendix D.1. \href{https://arxiv.org/abs/1205.0971}{Sen's ensemble calculation} \enspace/\enspace \repo{audit/ensemble-crosscheck.md}{Area-packet derivation} \enspace/\enspace \repo{audit/sen.md}{Coefficient comparison}.}

\chapterpage{tests}{What would decide the proposal?}{Research priorities}
\subhead{1. Fix the classical boundary problem}
Specify the data held fixed at both caps and the corner. Derive the two-form from the full action, including moving endpoints and joints. Repeat the endpoint variation on page 2 and verify that the new form is closed. Vary the proposed Hamiltonian on the same space and account for the flux. Record which area and shape variations survive this repair.

\subhead{2. Construct a clock at increasing resolution}
Start with one physical polarization and one angular harmonic. Solve the constraints before counting the radial modes. Keep the cap fraction, physical area, $G_N$, $\hbar$ and $\ell$ fixed. Refine radial and angular resolution separately, then together. Track the bottom of the clock spectrum and the probability $\|P_\ell\Omega\|^2$. A proof of a diverging minimum would exclude that clock prescription. Establishing a nonzero limiting overlap requires control of the quantum form and embeddings across cutoffs. A linearized model gives a mean squared shear that grows as the square of the radial cutoff. High angular modes may add little to the reading. Extending either result to interacting gravity requires control of all the constrained modes.

\subhead{3. Resolve finite rank in the physical space}
Search for a physical corner charge with continuous values that commutes with $P_\ell$. Disjoint spectral intervals of that charge could supply the states seen on page 4. Conversely, one sufficient positive result is a finite-dimensional projection $F$ with $\|(1-F)P_\ell\|<1$. Then $F$ is injective on the selected space, bounding its dimension. Either calculation must include edge excitations, the cap constraints and states reached by area-changing sources. A continuous classical parameter by itself leaves its quantum spectrum open.

\subhead{4. Test dynamics and closed sewing independently}
Construct regulated interacting evolution for a common set of states. Verify that different intermediate cuts give the same matrix element when the regulator is removed. Compare the resulting weak-field four-graviton amplitude with the Einstein tree amplitude. For the closed experiment, compare the unnormalized boost response for distinct exhaustive cuts against an independent uncut answer. A bound on the errors must hold as more states are included.

\subhead{5. Calculate the full packet measure}
Repeat the Schwarzschild calculation for the actual clock-filtered state, checking whether the trace selects the marginal horizon saddle over a flat round sphere. Evaluate the area integration and zero modes on joined replicas, keeping the same boundary conditions as in the state count. Recovering the spin-summed result and the changes when spin is restricted would test the ensemble identification directly.

\source{\repo{REFEREE.md}{Questions for a referee} \enspace/\enspace \repo{audit/regulator-program.md}{Regulator choices and pass criteria} \enspace/\enspace \repo{audit/finite-kinematics.md}{Physical corner-charge tests}.}


\chapterpage{observables}{What could a measurement distinguish?}{Four quantities to calculate or measure}
\subhead{An ellipse and a radar range}
Send a circular pencil of light against a vacuum plane wave and reflect its central ray back to a clock. On the wave axis, compare its round-trip radar range $L$ with the optical distance $D=\sqrt{dA/d\Omega}$ inferred from the spot area and launch angle. On a short constant-curvature segment, the signed spot ellipticity $e$ obeys
\[
     1-D/L=e^2/5+O(e^4).
\]
For an illustrative strain $10^{-21}$ at 100 Hz over 4 km, $|e|\simeq2.34\times10^{-26}$ and the predicted fractional distance deficit is $1.10\times10^{-52}$. The distance change of $4.39\times10^{-49}\,\mathrm m$ is about $10^{31}$ times smaller than the differential arm-length response $hL$ of an optimally oriented interferometer at that strain. The central echo time stays fixed in this geometry. This is a classical Einstein prediction for a screen and a clock. A comparison of quantizations needs a calculable change for the same apparatus. \repo{audit/observable-optics.md}{Optical derivation.}

\subhead{The spins of two falling masses}
Prepare each of two $10^{-14}\,\mathrm{kg}$ masses on a pair of paths labelled by spin. After a $2.5\,\mathrm{s}$ hold, recombine the paths and measure the two spins. For centres $450\,\mu\mathrm m$ apart and a $250\,\mu\mathrm m$ splitting per mass, coherent Newtonian gravity gives an ideal entangling phase of $-0.314$ radians. The predicted overlap probability with a specified entangled spin state is $0.578$. Every unentangled state gives at most $0.5$. A gravitational interpretation requires controls on electromagnetic forces, readout and initial correlations. A calculation from the four-dimensional vacuum proposal would first have to couple matter and specify the detectors. \repo{audit/observable-lab.md}{Spin observable and controls.}

\subhead{The fraction below half the allowed area}
In the proposed large Schwarzschild clock ensemble, form a ratio of physical traces with the same weight $e^{-s}$. Restrict the numerator to $s\le s_*/2$ and use all allowed areas in the denominator. With the paper's packet measure and its large-diamond clock-ensemble scaling, the predicted fraction is $2^{-257/90}\simeq0.13816$. An alternative assignment of the determinant to a density per unit area gives $2^{-151/45}\simeq0.09770$. An independent calculation must obtain the state weights and remove the regulator. An observational test would also need access to the prepared interior clock and area. \repo{audit/observable-entropy.md}{Trace ratio and assumptions.}

\subhead{Interference of two boundary frames}
If a reference apparatus could coherently boost one branch by rapidity $\eta$, its squared fringe visibility would be $V=|\langle e^{i\eta s/(2\pi)}\rangle|^2$. For two equally weighted areas separated by $\Delta s=4\pi^2$ (about 158 Planck areas), the first zero is at $\eta=0.5$. A uniform continuum of the same width gives its first zero at $\eta=1$. These example preparations distinguish area distributions. Adding more states at each area while preserving its total probability leaves the fringes unchanged. The controlled-boost apparatus remains a proposal. \repo{audit/observable-visibility.md}{Visibility curves.}

\source{\repo{OBSERVABLES.md}{Observable predictions, independent comparisons and runnable checks}.}

\chapterpage{reading}{Where the conclusions stand}{Reading and reproduction}
\subhead{What can already be used}
The constrained bulk theory on its gravitational branch gives classical Einstein dynamics. Null shear focuses the rays by the expected amount. A two-dimensional Jackiw--Teitelboim gravity example admits infinitely many orthogonal states with the same bounded clock distribution. The specified four-dimensional pair overlap sums exactly. Around a regular classical saddle, the stationary-phase equations and the one-loop response of observables averaged at fixed resolution have the expected form. Removing that averaging requires further short-distance control.

\subhead{What a completed state count would change}
If a nonempty four-dimensional clock window has finite rank, the area operator has only finitely many eigenvalues within that window, provided it preserves the selected space. Their spacing and multiplicities require further calculations. Extending the construction to a finite thermal trace in anti-de Sitter space also needs control of a closed sum. Inferring spatial topology from a two-boundary trace uses additional assumptions about factorization. A black-hole evaporation history would require dynamical amplitudes for radiation escaping the region.

An infinite-dimensional outcome would still leave room for useful entropy calculations. For example, some infinite-dimensional observable algebras admit a trace finite on the states of interest while the ordinary state count diverges. Its normalization fixes entropy up to a shared additive constant. The paper uses the ordinary trace. Deriving a different algebra for the diamond would require a separate calculation.

\subhead{Reproduce three checks}
From the \href{https://github.com/sprajs/quantum-gravity}{public review repository}, run the following with Python 3. They use its standard library.
\begin{center}\small
\begin{tabular}{@{}p{0.45\linewidth}p{0.5\linewidth}@{}}\toprule
Command & What to look for\\\midrule
\texttt{python3 checks/pair.py} & Compare the complex mode sum with its analytic tail bound.\\[4pt]
\texttt{python3 checks/corner\_clock.py} & At $\ell/u_0=1.1$ the accepted count grows from 6 to 94 across the printed matrix sizes.\\[4pt]
\texttt{python3 checks/05-2.py} & Compare exhaustive cut orders in the two-area model.\\\bottomrule
\end{tabular}\end{center}

\panel{Source and scope}{This nine-page handbook assesses the 83-page version posted on 25 November 2026. Both the manuscript and this technical review acknowledge AI assistance: Astra checks share a model family with the paper's collaborator, Fable's review comes from a separate family, and neither has had human specialist peer review. Readers can inspect every hypothesis, run the code and compare the exact calculations with the versioned preprint.}

\source{\href{https://arxiv.org/abs/2611.34924v1}{Versioned preprint} \enspace/\enspace \repo{EQUATIONS.md}{Printed equation finder} \enspace/\enspace \repo{audit/dependencies.md}{Dependencies among the claims} \enspace/\enspace \repo{docs/handbook.tex}{Editable handbook source}.}
\end{document}
